% =====================================================================
%  Shared preamble -- owned by LEAD.
%  Section authors: do not add packages here; request them in NOTES-<owner>.md.
%  Ported from Paper 2 (Shock and Contagion), itself a port of Paper 1 (The
%  Dark Figure), with Paper 2's later fixes (run-in heads that never double
%  their period; section titles that cannot be stranded from their rule), so
%  that the four papers read as a series.
% =====================================================================

% ---- fonts & encoding (amsthm MUST precede newtxmath) ----------------
\usepackage[T1]{fontenc}
\usepackage[utf8]{inputenc}
\usepackage{amsmath,mathtools,amsthm}
\IfFileExists{newtxtext.sty}{%
  \usepackage{newtxtext}\usepackage[varvw]{newtxmath}%
}{%
  \usepackage{mathptmx}\usepackage{amssymb}%
}
% newtx's TS1 (text-companion) metrics, ts1-qtm*, are not installed here; map
% TS1 onto Adobe Times (psnfss) so \textbullet, \textdegree, etc. still work.
% Without this, a plain itemize (\textbullet) aborts the run.
\DeclareFontFamily{TS1}{ntxtlf}{}
\DeclareFontShape{TS1}{ntxtlf}{m}{n}{<->ssub * ptm/m/n}{}
\DeclareFontShape{TS1}{ntxtlf}{m}{it}{<->ssub * ptm/m/it}{}
\DeclareFontShape{TS1}{ntxtlf}{m}{sl}{<->ssub * ptm/m/it}{}
\DeclareFontShape{TS1}{ntxtlf}{m}{sc}{<->ssub * ptm/m/n}{}
\DeclareFontShape{TS1}{ntxtlf}{b}{n}{<->ssub * ptm/b/n}{}
\DeclareFontShape{TS1}{ntxtlf}{b}{it}{<->ssub * ptm/b/it}{}
\DeclareFontShape{TS1}{ntxtlf}{b}{sl}{<->ssub * ptm/b/it}{}
\DeclareFontShape{TS1}{ntxtlf}{b}{sc}{<->ssub * ptm/b/n}{}
\DeclareFontShape{TS1}{ntxtlf}{bx}{n}{<->ssub * ptm/b/n}{}
\DeclareFontShape{TS1}{ntxtlf}{bx}{it}{<->ssub * ptm/b/it}{}
\DeclareFontShape{TS1}{ntxtlf}{bx}{sl}{<->ssub * ptm/b/it}{}
\DeclareFontShape{TS1}{ntxtlf}{bx}{sc}{<->ssub * ptm/b/n}{}
\usepackage[letterpaper,margin=1.05in,headheight=14pt,headsep=18pt,footskip=30pt]{geometry}
\usepackage{etoolbox}
\usepackage{thmtools}
\usepackage{bm}
\usepackage[dvipsnames,table]{xcolor}

% ---- house palette (deep navy + one burnt-orange accent) -------------
%  Same values as Papers 1-3. In this paper navy carries structure and the
%  algorithmic policy; light navy the reference policy and neutral content;
%  mid navy borderline content; the accent marks antisemitic content and what
%  is fixed in advance (the stopping boundary, the registered margin).
\definecolor{LabNavy}{HTML}{14284B}
\definecolor{LabNavyMid}{HTML}{3D5A80}
\definecolor{LabNavyLight}{HTML}{8DA1BE}
\definecolor{LabAccent}{HTML}{B8460E}   % print-safe variant of the lab orange
\definecolor{LabInk}{HTML}{1B1D22}
\definecolor{LabGrey}{HTML}{6B7079}
\definecolor{LabRule}{HTML}{C8CEDA}
\definecolor{LabTint}{HTML}{F2F4F8}
\definecolor{LabSand}{HTML}{FBF7F3}

% ---- microtypography --------------------------------------------------
\usepackage[final,protrusion=true,expansion=true,tracking=true,
            factor=1100,stretch=12,shrink=12]{microtype}
\SetTracking{encoding={*},shape=sc}{45}
\linespread{1.06}
\setlength{\emergencystretch}{1.5em}
\clubpenalty=3000
\widowpenalty=3000
\displaywidowpenalty=10000

\usepackage{booktabs,array,tabularx,multirow}
\usepackage{placeins}     % \FloatBarrier at the end of Section 4
\usepackage{graphicx}
\usepackage{caption,subcaption}
\usepackage[shortlabels]{enumitem}
\usepackage{algorithm}
\usepackage[noend]{algpseudocode}
\usepackage{tikz}
\usetikzlibrary{arrows.meta,positioning,calc,shapes.geometric,shapes.misc,fit,
  backgrounds,decorations.pathreplacing,matrix,fadings}
\usepackage{tcolorbox}
\tcbuselibrary{skins,breakable}
\usepackage{titlesec}
\usepackage[authoryear,round,sort]{natbib}
\usepackage{fancyhdr}
\usepackage{xurl}        % long URLs/DOIs in the bibliography may break anywhere
\usepackage[colorlinks=true,linkcolor=LabNavy,citecolor=LabNavy,
            urlcolor=LabNavyMid,bookmarksnumbered=true,pdfdisplaydoctitle=true,
            hypertexnames=false]{hyperref}
\usepackage[capitalise,noabbrev,nameinlink]{cleveref}
% Unique hyperref anchors for algorithm lines (several algorithms number their
% lines from 1; without this hyperref reports duplicate destinations).
\makeatletter\def\theHALG@line{\arabic{algorithm}.\arabic{ALG@line}}\makeatother

% ---- bibliography -----------------------------------------------------
\urlstyle{same}
\setlength{\bibsep}{2.5pt plus 1pt}
\renewcommand{\bibfont}{\small}
\setlength{\bibhang}{1.2em}

% ---- captions ---------------------------------------------------------
\DeclareCaptionLabelFormat{labnavy}{\textcolor{LabNavy}{\bfseries #1~#2}}
\captionsetup{font=small,labelformat=labnavy,labelsep=period,
  justification=justified,singlelinecheck=true,skip=6pt}
\captionsetup[table]{position=top,skip=4pt}
\captionsetup[sub]{font=small,labelformat=parens,labelfont=bf}
\captionsetup[algorithm]{font=normalsize,labelformat=labnavy,labelsep=space}

\setlength{\parskip}{0pt}
\setlist{itemsep=2pt,topsep=4pt}
\setlist[itemize,1]{label={\color{LabNavy}\small$\bullet$}}
\setlist[itemize,2]{label={\color{LabNavyMid}\textendash}}
\setlist[itemize,3]{label={\color{LabNavyMid}\small$\circ$}}
\renewcommand{\arraystretch}{1.12}
\setlength{\aboverulesep}{0.5ex}
\setlength{\belowrulesep}{0.6ex}

% ---- section headings ---------------------------------------------------
%  Navy titles, accent-coloured numbers; a hairline under each section.
%  \nobreak keeps the hairline (and the text) with the title: without it a
%  page can break between a section title and its rule.
\titleformat{\section}[hang]
  {\normalfont\Large\bfseries\color{LabNavy}}
  {\textcolor{LabAccent}{\thesection}}{0.65em}{}
  [{\nobreak\vspace{2pt}\nobreak\color{LabRule}\nobreak\titlerule[0.5pt]}]
\titlespacing*{\section}{0pt}{3.0ex plus 1ex minus .3ex}{1.6ex plus .3ex}
\titleformat{\subsection}[hang]
  {\normalfont\large\bfseries\color{LabNavy}}
  {\textcolor{LabNavyMid}{\thesubsection}}{0.6em}{}
\titlespacing*{\subsection}{0pt}{2.4ex plus .8ex minus .3ex}{1.0ex plus .2ex}
\titleformat{\subsubsection}[hang]
  {\normalfont\normalsize\bfseries\itshape\color{LabNavy}}
  {\thesubsubsection}{0.55em}{}
\titlespacing*{\subsubsection}{0pt}{2.0ex plus .6ex minus .2ex}{0.7ex plus .2ex}
% Run-in heads end with a period unless the author already typed one
% (\paragraph{Streams} and \paragraph{Streams.} both give "Streams.").
\makeatletter
\newcommand{\labparatitle}[1]{#1\@addpunct{.}}
\makeatother
\titleformat{\paragraph}[runin]
  {\normalfont\normalsize\bfseries\color{LabNavy}}{\theparagraph}{0.5em}{\labparatitle}
\titlespacing*{\paragraph}{0pt}{1.4ex plus .5ex minus .2ex}{0.6em}
% Appendix sections read "Appendix A  Title".
\appto\appendix{%
  \titleformat{\section}[hang]
    {\normalfont\Large\bfseries\color{LabNavy}}
    {\textcolor{LabAccent}{\appendixname~\thesection}}{0.65em}{}
    [{\nobreak\vspace{2pt}\nobreak\color{LabRule}\nobreak\titlerule[0.5pt]}]%
}

% ---- running heads --------------------------------------------------------
\pagestyle{fancy}
\fancyhf{}
\fancyhead[L]{\footnotesize\scshape\color{LabGrey}Supply or Demand?}
\fancyhead[R]{\footnotesize\color{LabGrey}Kornstein Lab for Antisemitism Research
  \enspace\textcolor{LabRule}{|}\enspace\textcolor{LabNavy}{\bfseries\thepage}}
\renewcommand{\headrulewidth}{0.4pt}
\renewcommand{\headrule}{{\color{LabRule}%
  \hrule width\headwidth height\headrulewidth \vskip-\headrulewidth}}
\fancypagestyle{plain}{%
  \fancyhf{}%
  \fancyfoot[C]{\footnotesize\color{LabGrey}\thepage}%
  \renewcommand{\headrulewidth}{0pt}%
  \renewcommand{\headrule}{}%
}

% ---- footnotes ----------------------------------------------------------------
\renewcommand{\footnoterule}{{\color{LabRule}\kern-3pt\hrule width 0.28\columnwidth height 0.4pt\kern 2.6pt}}

% ---- abstract & keywords ---------------------------------------------------
\renewenvironment{abstract}{%
  \par\vspace{2pt}%
  \begin{list}{}{\leftmargin=0pt\rightmargin=0pt\topsep=0pt\partopsep=0pt}%
  \item[]{\footnotesize\scshape\bfseries\color{LabAccent}Abstract}\par\vspace{3pt}%
  \small\noindent\ignorespaces}%
 {\par\end{list}}
\newcommand{\keywords}[1]{\par\vspace{5pt}\noindent{\small
  {\scshape\bfseries\color{LabNavy}Keywords}\enspace #1}\par}

% ---- theorem environments (one shared counter, numbered within section) ---
\newtheoremstyle{labplain}{4pt}{4pt}{\itshape}{}{\bfseries\color{LabNavy}}{.}{.5em}
  {\thmname{#1}\thmnumber{ #2}\thmnote{\normalfont\color{LabNavy}\ (#3)}}
\newtheoremstyle{labdef}{4pt}{4pt}{\normalfont}{}{\bfseries\color{LabNavy}}{.}{.5em}
  {\thmname{#1}\thmnumber{ #2}\thmnote{\normalfont\color{LabNavy}\ (#3)}}
\newtheoremstyle{labremark}{6pt}{6pt}{\normalfont}{}{\itshape\color{LabNavy}}{.}{.5em}
  {\thmname{#1}\thmnumber{ #2}\thmnote{\ (#3)}}

\theoremstyle{labplain}
\newtheorem{theorem}{Theorem}[section]
\newtheorem{proposition}[theorem]{Proposition}
\newtheorem{lemma}[theorem]{Lemma}
\newtheorem{corollary}[theorem]{Corollary}
\theoremstyle{labdef}
\newtheorem{definition}[theorem]{Definition}
\newtheorem{assumption}{Assumption}
\renewcommand{\theassumption}{A\arabic{assumption}}
\theoremstyle{labremark}
\newtheorem{example}[theorem]{Example}
\newtheorem{remark}[theorem]{Remark}
\crefname{assumption}{Assumption}{Assumptions}
\crefname{algorithm}{Algorithm}{Algorithms}
\crefname{appendix}{Appendix}{Appendices}
\numberwithin{equation}{section}

% Subtle left rule for formal results: navy for theorem-like statements,
% light navy for definitions and assumptions. Breakable across pages.
\tcbset{labbar/.style={blanker,breakable,
  left=9pt,right=0pt,top=1pt,bottom=1pt,
  before skip=8pt plus 2pt,after skip=8pt plus 2pt,
  borderline west={1.1pt}{0pt}{#1}}}
% thmtools hooks (not \tcolorboxenvironment) so that restatable/restated
% results receive the same rule.
\newcommand{\labbarthm}[2]{%
  \addtotheorempreheadhook[#1]{\tcolorbox[labbar=#2]}%
  \addtotheorempostfoothook[#1]{\endtcolorbox}}
\labbarthm{theorem}{LabNavy}
\labbarthm{proposition}{LabNavy}
\labbarthm{corollary}{LabNavy}
\labbarthm{lemma}{LabNavyMid}
\labbarthm{definition}{LabNavyLight}
\labbarthm{assumption}{LabAccent!70}

% Algorithms: navy rules and caption label.
\makeatletter
\renewcommand\fs@ruled{\def\@fs@cfont{\bfseries\color{LabNavy}}\let\@fs@capt\floatc@ruled
  \def\@fs@pre{{\color{LabNavy}\hrule height.8pt depth0pt}\kern2pt}%
  \def\@fs@post{\kern2pt{\color{LabNavy}\hrule height.5pt}\relax}%
  \def\@fs@mid{\kern2pt{\color{LabNavy!60}\hrule height.4pt}\kern2pt}%
  \let\@fs@iftopcapt\iftrue}
\makeatother

\renewcommand{\qedsymbol}{\textcolor{LabNavyMid}{\ensuremath{\square}}}
\makeatletter
\renewenvironment{proof}[1][\proofname]{\par
  \pushQED{\qed}\normalfont\topsep6\p@\@plus6\p@\relax
  \trivlist\item[\hskip\labelsep{\itshape\color{LabNavy}#1}\@addpunct{.}]\ignorespaces
}{\popQED\endtrivlist\@endpefalse}
\makeatother

% ---- notation macros (Paper 1; kept because section authors may use them) --
%  Paper 4 notation (AR, AA, Se, Sp, P_alg, ...) is written out in full in the
%  sections (e.g. \mathrm{AR}); section-local helpers such as \osc and \dTV
%  are declared with \providecommand where they are used.
\newcommand{\E}{\mathbb{E}}
\newcommand{\dW}{d_{\mathrm{W}}}            % Wasserstein-1 distance
\newcommand{\erg}{\operatorname{erg}}       % ergodicity coefficient
\renewcommand{\Pr}{\mathbb{P}}
\newcommand{\Var}{\operatorname{Var}}
\newcommand{\Cov}{\operatorname{Cov}}
\newcommand{\ind}[1]{\mathbf{1}\!\left\{#1\right\}}
\newcommand{\Kset}{[K]}
\newcommand{\Hset}{\mathcal{H}}
\newcommand{\Hobs}{\mathcal{H}_{+}}          % nonzero histories
\newcommand{\zero}{\mathbf{0}}
\newcommand{\h}{\bm{h}}
\newcommand{\nzero}{n_{\zero}}
\newcommand{\Nhat}{\widehat{N}}
\newcommand{\Link}{\Lambda}
\newcommand{\Linkhat}{\widehat{\Lambda}}
\newcommand{\Lop}{\mathcal{T}}               % linkage operator T(delta, eps)
\newcommand{\dirichlet}{\operatorname{Dir}}
\newcommand{\DP}{\operatorname{DP}}
\newcommand{\Bern}{\operatorname{Bernoulli}}
\newcommand{\Pois}{\operatorname{Poisson}}
\newcommand{\Mult}{\operatorname{Multinomial}}
\newcommand{\Bin}{\operatorname{Binomial}}
\newcommand{\supp}{\operatorname{supp}}
\newcommand{\logit}{\operatorname{logit}}
\newcommand{\argmin}{\operatorname*{arg\,min}}
\newcommand{\argmax}{\operatorname*{arg\,max}}
\newcommand{\rvic}{\rho_{\mathrm{vic}}}
\newcommand{\rpol}{\rho_{\mathrm{pol}}}
\newcommand{\rcls}{\rho_{\mathrm{cls}}}
\newcommand{\rsub}{\rho_{\mathrm{sub}}}
\newcommand{\nFBI}{n_{\mathrm{FBI}}}
\newcommand{\factcheck}[1]{\textcolor{red}{[FACT-CHECK: #1]}}
\newcommand{\todo}[1]{\textcolor{orange}{[TODO: #1]}}

